% УДИРТГАЛ — БҮЛЭГ 1 (дугааргүй бүлэг)
% Эх: diploma_0.2v.docx, БҮЛЭГ 1. УДИРТГАЛ

\phantomsection
\addchaptertocentry{Удиртгал}

%-------------------------------------------------------------------------------
\section*{Удиртгал}

Интернэтийн хэрэглэгчдийн тоо дэлхий даяар урьд өмнө байгаагүй хурдацтай
нэмэгдэж байгаа бөгөөд цахим мэдээллийн орчин нийгмийн харилцааны хэлбэрийг
бүрэн өөрчилсөн. DataReportal (2025)-ийн мэдээллийн дагуу Монгол Улсын
интернэт хэрэглэгчдийн тоо нийт хүн амын 83 гаруй хувийг эзэлж байгаа бөгөөд
Facebook, X (Twitter), TikTok зэрэг платформуудын хэрэглэгчид жил бүр хурдацтай
өсч байна. Цахим орчинд тархаж буй мэдээлэл, харилцааны хэмжээ нэмэгдэхийн
хэрээр хэрэглэгчид, ялангуяа өсвөр насныхан, эрүүл бус, хортой агуулгад
өртөх эрсдэл ч нэгэн зэрэг нэмэгдэж байна.

Монголын цахим орчинд --- тухайлбал news.mn зэрэг мэдээний сайтын сэтгэгдлийн
хэсэг, Facebook-ийн нийтийн хуудас болон бүлгүүд, X (Twitter)-ийн нийтийн
нийтлэлүүдэд --- кибер дарамт, үзэн ядалтын хэл, хорлолтой дезинформаци, хэт
аггрессив хандлага зэрэг олон хэлбэрийн сөрөг агуулга тархаж байна. Олон
улсын судалгаанууд \cite{kowalski2014,vogels2022} кибер дарамт нь өсвөр насны
хүүхдийн сэтгэл зүйн тогтвортой байдал, нийгмийн харилцаа, академик
гүйцэтгэлд ноцтой сөргөөр нөлөөлдгийг нотолсон. Монгол Улсад энэ асуудалд
хариулах техникийн шийдэл, тухайлбал автомат контент шүүлтийн систем,
хөгжлийн түвшин хангалтгүй байгааг тэмдэглэх нь зүйтэй.

Гэхдээ аливаа сөрөг агуулгыг цогцоор нь хориглох нь нийгмийн хувьд буруу
шийдэл болно. Бүтээлч шүүмжлэл, асуудалд суурилсан шинжилгээ, бодлогыг
эсэргүүцэх тайлбар зэрэг нь ардчилсан нийгэм, иргэний соёлын чухал суурь
болдог. ``Positive-only'' орчин нь нийгмийн бодит байдлаас тасарсан хиймэл
дискурсийг бий болгох аюултай бөгөөд үзэл бодлоо чөлөөтэй илэрхийлэх нь
үндсэн хүний эрхэд хамаарна. Тиймд автомат шүүлтийн систем нь зөвхөн бүх
сөрөг агуулгыг устгах зорилготой байж болохгүй, харин хоёр эсрэг туйлын
агуулгыг --- нэг нь хориглох, нөгөөг нь хамгаалах --- нарийн ялгаж чаддаг
байх шаардлагатай.

Монгол хэлний тусгайлсан NLP судалгааны хүрээнд Toxic болон Constructive
агуулгыг ялгах систем одоогоор боловсруулагдаагүй байгаа нь судалгааны
цоорхойг бүрдүүлж байна. Өмнөх ажлуудын дийлэнх нь гурван ангилалт
(Positive / Negative / Neutral) сентимент шинжилгээгээр хязгаарлагдаж,
``Negative'' ангилалын дотоодод байх хоёр эсрэг шинж чанарыг --- устгах
шаардлагатай хортой агуулга ба хамгаалах шаардлагатай Бүтээлч шүүмжлэл ---
ялган тодорхойлоогүй байна. Энэхүү судалгааны ажил тэрхүү цоорхойг нөхөхийг
зорьж буй бөгөөд нийгмийн болон техникийн ач холбогдлоороо тодорхой хувь
нэмэр оруулна.

%-------------------------------------------------------------------------------
\section*{Судалгааны зорилго}

Монгол хэл дээрх цахим орчны агуулгыг олон ангилалтаар шүүх, тухайлбал Toxic
Negative болон Constructive Criticism-ийг нарийн ялгах чадвартай хиймэл оюунд
суурилсан NLP системийг боловсруулахын тулд шатлалт (hierarchical) болон нэг
шатлалт (flat) архитектуруудыг харьцуулан шинжилж, хамгийн өндөр гүйцэтгэлтэй
оновчтой шийдлийг тодорхойлоход энэхүү судалгааны гол зорилго оршино.

%-------------------------------------------------------------------------------
\section*{Судалгааны зорилтууд}

Дипломын ажлын хүрээнд дараах зорилтуудыг дараалсан байдлаар хэрэгжүүлнэ.

\begin{enumerate}
  \item news.mn болон Facebook нийтийн хуудас, бүлгүүдээс Монгол хэл дээрх
        сэтгэгдэл болон нийтлэлүүдийг цуглуулах, цэвэрлэх, стандартжуулах
        ажлыг гүйцэтгэнэ.

  \item Цуглуулсан өгөгдлийг Positive, Neutral, Toxic Negative, Constructive
        Criticism гэсэн дөрвөн ангилалд 3 хүний бүрэлдэхүүнтэйгээр гар аргаар
        шошголох (manual annotation) бөгөөд аннотаторуудын хоорондын нийцлийг
        статистик аргаар хэмжинэ.

  \item Монгол хэлний agglutinative онцлогт тохирсон урьдчилсан боловсруулалтын
        (preprocessing) pipeline --- токенчлол, нормчлол --- боловсруулна.

  \item Хоёр өөр архитектурын туршилтыг зохион байгуулна: (1) Хоёр шатлалт
        каскад загвар болон (2) Focal Loss бүхий нэг шатлалт (Flat) MN-BERT
        загвар.

  \item MN-BERT загварыг fine-tuning хийж Logistic Regression, Naive Bayes,
        SVM зэрэг суурь загваруудтай харьцуулан статистик аргаар үнэлнэ.

  \item Алдааны дамжуулалт (Error Propagation) болон ангиллын хил хязгаарын
        дүн шинжилгээнд үндэслэн хамгийн оновчтой архитектурыг сонгоно.

  \item Сонгосон оновчтой загварт суурилсан Streamlit/Gradio хэрэглэгчийн
        интерфейстэй прототип системийг хэрэгжүүлж, туршилтын орчинд үнэлнэ.
\end{enumerate}

%-------------------------------------------------------------------------------
\section*{Судалгааны хамрах хүрээ}

Энэхүү судалгаа нь дараах хэмжээст орон зайд явагдана. Хэлний хувьд зөвхөн
Монгол кирилл бичгийн текст хамрагдана. Латин үсгийн бичвэр, монгол-англи
холимог (code-switched) текст, романчилсан Монгол бичлэг (``bn'', ``yu ve'',
``zgr'' зэрэг), болон бусад кирилл бус бичиглэлийн хэлбэрүүд датасетаас
тусгайлан хасагдана. Энэхүү шийдэл нь MN-BERT-ийн кирилл корпус дээр суурилсан
токенчлолын тогтвортой байдлыг хангаж, аннотацийн стандартыг нэгтгэхэд
шаардлагатай.

Платформын хувьд news.mn сэтгэгдлийн хэсэг болон Facebook нийтийн хуудас,
нийтийн нийтлэлүүд хамрагдана. Хугацааны хувьд 2024--2025 оны нийтлэлүүдийг
ашиглана. Контентийн хувьд зөвхөн текстэн өгөгдлийг анализад оруулах бөгөөд
зураг, видео, эможи зэргийг анхдагч загварт оруулахгүй --- эдгээрийг ирээдүйн
судалгааны чиглэл болгон тодорхойлно. Системийн хувьд прототип загвар
боловсруулахаар тооцох бөгөөд production орчинд байршуулах нь энэхүү ажлын
хүрээнд багтахгүй.

%-------------------------------------------------------------------------------
\section*{Бүлгийн дүгнэлт}

Монгол хэлний цахим орчинд хортой болон Бүтээлч (Constructive) агуулгыг
автоматаар ялган шүүх асуудал нь зөвхөн техникийн биш, нийгэм, иргэний болон
хүний эрхийн чухал ач холбогдолтой. Өмнөх судалгааны ажлуудад энэ ялгалт
хийгдэж байгаагүй нь тодорхой судалгааны цоорхойг бий болгосон. Энэхүү ажил
тэрхүү цоорхойг нөхөж, Монгол хэлний NLP-ийн хөгжилд хувь нэмэр оруулна.
Цаашдын бүлгүүдэд онолын суурь, арга зүй, хэрэгжүүлэлт, туршилтын үр дүнг
дэлгэрэнгүй танилцуулна.
